\chapter{Installing TandemLLM}\label{ch:install-engine}

TandemLLM has two parts to install. The engine runs the models and needs a CUDA GPU; its test board is the DGX Spark (Section~\ref{sec:test-env}). The dashboard and its proxy need only Python, run on any machine, and also run inside the engine's own server. Most setups use both. Find your setup in Table~\ref{tab:setups}, then follow the sections it names.

\section{Setups, and which ones work}\label{sec:setups}\index{setups}

The parts combine. One engine can sit behind one address together with other engines and cloud APIs, and one dashboard watches all of them. Table~\ref{tab:setups} says which combinations work today, and how far each one is tested.

\begin{table}[htbp]
\centering\small
\caption{Setups and their state in this release. \emph{Live}: run on the test board in the 0.4.0 series. \emph{Built}: the code supports it and its tests pass, but no live run is recorded.}\label{tab:setups}
\begin{tabularx}{\linewidth}{@{}L{.27\linewidth}YL{.13\linewidth}@{}}
\toprule
\hd{setup} & \hd{how} & \hd{state} \\
\midrule
the engine alone & \code{install.sh} on the board (this chapter) & live \\
the engine and cloud APIs, one address & \code{tandem-dash} on the board or another machine, with the engine and the cloud API as backends (Chapter~\ref{ch:install-dash}) & live, with DeepSeek \\
the same, from the engine's own server & \code{QSE_BACKENDS}: the engine relays to the cloud API, and its conductor can run in front of it (Section~\ref{sec:engine-backends}) & built \\
the engine and other engines on other machines & either of the two above, the other engines as \code{vllm}, \code{sglang} or \code{generic} backends (Chapter~\ref{ch:other-engines}) & built; vLLM and SGLang live behind the proxy on their own \\
the dashboard and proxy alone & \code{pip install} or docker on any machine, in front of the backends you run (Chapter~\ref{ch:install-dash}) & live \\
the engine and a second GPU engine on one larger machine & each engine on its own GPU, \code{tandem-dash launch} for the second & built \\
the engine and a second GPU engine on one DGX Spark & the engine plans its memory as if it owns the board (up to 77~GiB during a prefill, plus its small models); a board out of memory stops answering until it is power cycled & not supported \\
\bottomrule
\end{tabularx}
\end{table}

\begin{figure}[htbp]
\centering
% Which install path: a decision tree over the three paths of Table tab:install-paths.
\begin{tikzpicture}[node distance=6mm and 8mm,
    q/.style={draw=slate, fill=paper, shape=rounded rectangle, font=\sffamily\small,
      align=center, inner sep=4pt, minimum height=9mm},
    path/.style={draw=cobalt, fill=cobaltlight, rounded corners=2pt, text width=50mm, align=left,
      font=\sffamily\small, inner sep=5pt}]
  \node[q] (spark) {a machine with a CUDA GPU?\\{\scriptsize tested: DGX Spark}};
  \node[q, below right=8mm and 6mm of spark] (others) {other engines or\\cloud APIs too?};
  \node[q, below left=8mm and 6mm of spark] (py) {Python 3.11 or newer,\\or docker?};
  \node[path, below=9mm of others] (a) {\textbf{Path A: the engine}\\[2pt]
    {\normalfont\footnotesize the engine with its models, the conductor and the dashboard at \texttt{:8000}}\\[3pt]
    {\ttfamily\scriptsize curl -fsSL .../install.sh | bash}\\[2pt]
    {\footnotesize\color{slate} Chapter~\ref{ch:install-engine}}};
  \node[path, below=9mm of py] (b) {\textbf{Path B: the dashboard and proxy}\\[2pt]
    {\normalfont\footnotesize one address, one ledger and one dashboard for backends you already run; no GPU, no torch. In docker: \texttt{ops/Dockerfile.dash}}\\[3pt]
    {\ttfamily\scriptsize pip install '.[host]' \&\& tandem-dash}\\[2pt]
    {\footnotesize\color{slate} Chapter~\ref{ch:install-dash}}};
  \node[path, draw=signal, fill=signallight, text width=70mm, below=10mm of $(a.south)!0.5!(b.south)$] (c) {\textbf{Path C: several engines and cloud APIs behind one address}\\[2pt]
    {\normalfont\footnotesize A or B, plus the backends. The engine's own server relays (\texttt{QSE\_BACKENDS}), or \texttt{tandem-dash} lists the engine as one backend among others. Chapter~\ref{ch:other-engines}}};
  \draw[flow] (spark) -- node[right, lbl, pos=.55] {yes} (others);
  \draw[flow] (spark) -- node[left, lbl, pos=.55] {no} (py);
  \draw[flow] (others) -- node[right, lbl] {no} (a);
  \draw[flow] (py) -- node[left, lbl] {yes} (b);
  \draw[flow] (others.east) -- ++(8mm,0) |- node[pos=.25, right, lbl] {yes} (c.east);
  \draw[dflow] (a.south) -- node[right, lbl] {add backends} (a.south |- c.north);
  \draw[dflow] (b.south) -- node[left, lbl] {add backends} (b.south |- c.north);
\end{tikzpicture}

\caption[Which install path]{Which install path. A machine with the GPU gets the engine (path A); any other machine gets the dashboard and proxy (path B); either grows into path C by adding backends. The engine's own server can relay to other backends, and a standalone proxy can list the engine as one backend among others.}\label{fig:install-paths}
\end{figure}

\section{What the engine needs}\index{requirements}

The engine is tested on one NVIDIA DGX Spark, a GB10 GPU (\code{sm_121}) with 128~GB of memory that the CPU and the GPU share, running Linux on aarch64. On another CUDA GPU the installer turns off the one kernel built for the GB10, and the engine runs best effort. The installer checks four things before it changes anything: an NVIDIA GPU and driver, the CUDA compiler (one kernel builds on its first use), about 70~GB of free disk, and 64~GiB of free memory at the start. With \code{--setup} it downloads only Wren's model, so it asks for about 5~GiB of disk beside the Python environment and 8~GiB of free memory; \code{use-profile.sh} checks the disk and the memory a profile needs before it downloads or starts it.

The software is pinned. The engine imports PyTorch 2.13 built for CUDA 13.0, torchvision 0.28 (the CUDA 13.0 build that matches that torch), Triton 3.7, transformers 5.12, Pillow 11.3, safetensors and numpy, and nothing else. Python 3.11 is the tested version; under 3.12 one grammar test fails. The installer's lock lists every package with its version, for aarch64 and x86\_64.

Two of those packages are there for images. Pillow decodes a request's images, and torchvision is the backend of the image processor that the checkpoint ships. The engine pins that backend (\code{--image-backend torchvision}) and refuses to start without it, because transformers would otherwise fall back to its Pillow backend, whose patches differ by up to two 8-bit levels after a resize (\measured, 0.4~\% of the values of one resized image). Every cache keys on those patches, so the backend has to be fixed. \code{--image-backend pil} is the explicit choice without torchvision, and \code{--vision off} serves text only.

Run one engine per board\index{one engine per board}. Two engines loading side by side run the board out of memory, and a DGX Spark in that state stops answering until it is power cycled. Section~\ref{sec:one-engine} explains how the scripts enforce it.

\section{A new box: setup mode, then the Setup Guide}\label{sec:new-box}\index{setup mode}\index{use-profile.sh@\texttt{use-profile.sh}}

Since 0.8.1 a new DGX Spark needs two commands and no Hugging Face token. Every file a profile reads is public on Hugging Face: our own packs, the vendors' releases, and turboderp's Flash-Next packs (Chapter~\ref{ch:flashnext}).

\begin{shell}
curl -fsSLO https://raw.githubusercontent.com/0xBakeer/TandemLLM/v0.13.0/install.sh
bash install.sh --setup --yes
\end{shell}

The first line fetches the installer from the release tag. With \code{--setup} the installer puts the code and the Python environment in place and downloads only Wren's model (Qwen3.5-2B, 4.3~GB), the setup guide's assistant. It writes the configuration with no weights and starts setup mode on \code{127.0.0.1:8000}: the dashboard and its \texttt{\#/setup} page, Wren, and no model loaded.

Then open the Setup Guide in a browser and pick a model. The guide shows each model with its weight packs, checks that the choice fits in memory, downloads it and starts the engine; it walks through the model, the weights, the tools and the parallel slots in eight steps. \code{docs/setup.md} describes the page.

\begin{itemize}
\item On the box itself, open \texttt{http://127.0.0.1:8000/dashboard/\#/setup}.
\item From another computer, open an SSH tunnel to the box first, \code{ssh -L 8000:127.0.0.1:8000 <user>@<box>}, and open the same address on that computer.
\item On your network, install with \code{--host 0.0.0.0} and open \texttt{http://<box>:8000/dashboard/\#/setup} from any computer that reaches the box. Until an admin password is set, anyone who reaches the port can run the setup.
\end{itemize}

The installer prints the addresses when it finishes, and while the box is in setup mode the dashboard opens on the Setup Guide by itself.

\subsection{Advanced: from the command line}\label{sec:new-box-cli}

The Setup Guide runs \code{ops/use-profile.sh} for you. To pick the weights from a shell instead, run the copy in \code{~/TandemLLM/run/ops/}, next to the \code{serve.env} the installer wrote:

\begin{shell}
bash ~/TandemLLM/run/ops/use-profile.sh --list
bash ~/TandemLLM/run/ops/use-profile.sh <profile> --dry-run
bash ~/TandemLLM/run/ops/use-profile.sh <profile>
\end{shell}

\code{--list} prints every profile and marks the active one. With \code{--dry-run} the script prints every file's URL and size and changes nothing. Without it, the script downloads what the profile needs and checks every weight file against the sha256 the hub lists. It then leaves setup mode, starts the engine on port 8000 and sends one smoke request. Since 0.8.2, \code{community/<id>} runs a pack someone else published (Chapter~\ref{ch:community}).

\begin{table}[htbp]
\centering\small
\caption{What two of the profiles download on a new box. Free disk: the download plus the 60~GB the script keeps free after it.}\label{tab:new-box}
\begin{tabularx}{\linewidth}{@{}L{.17\linewidth}YL{.14\linewidth}L{.1\linewidth}@{}}
\toprule
\hd{profile} & \hd{downloads} & \hd{size} & \hd{free disk} \\
\midrule
\code{flashnext-3.05} & the pack's branch of \code{turboderp/Qwen3.8-Flash-Next-exl3} & 85.1~GB & 145~GB \\
\code{exl3-4.5} & \code{0xBakeer/TandemLLM-Qwen3.8-27B-EXL3-4.5}; the K3 drafters and the K4 head (\code{...-EXL3-drafters}); the BF16 drafters (\code{...-DFlash2-b8}, \code{-b16}, into \code{~/TandemLLM/src/train}); the base checkpoint \code{Qwen/Qwen3.8-27B-FP8}, without \code{mtp.safetensors} & 18.4 + 1.9 + 7.7 + 30.4~GB & 118~GB \\
\bottomrule
\end{tabularx}
\end{table}

A folder appears only when all its files are there. If a download stops, run the same command again: it keeps the files that finished and continues. The engine then needs about 80~GiB of free memory for a 27B profile and 70~GiB for \code{flashnext-3.05}. If the profile does not come up, the script brings setup mode back.

The installer checks out the tag \code{v0.13.0} from the repository \code{--repo} names, the public one by default. Where the tag is missing, \code{curl} fails with a 404, and the installer stops at its clone step with \enquote{no tag, branch or commit 'v0.13.0'}. \code{--ref} picks another release, and \code{--repo} another repository.

\section{Updating a box in setup mode}\label{sec:update-setup}\index{upgrade}

A box with a running profile counts as configured, and \code{--setup} alone refuses to touch it. To update, add \code{--force-setup}: the installer stops the engine, checks out the new tag and comes back in setup mode. Then switch to the profile again; its files are on disk already, so nothing downloads twice.

\begin{shell}
bash ~/TandemLLM/src/install.sh --setup --force-setup --yes --ref <new tag>
bash ~/TandemLLM/run/ops/use-profile.sh flashnext-3.05
\end{shell}

\section{The installer}\index{install.sh}

Without \code{--setup}, one command sets up the NVFP4 profile, StairCut included:

\begin{shell}
curl -fsSL https://raw.githubusercontent.com/0xBakeer/TandemLLM/v0.13.0/install.sh | bash
\end{shell}

The installer works in seven steps, and each step is skipped when it is already done, so running it again updates an install and resumes a download that stopped.

\begin{enumerate}
\item Checks. The GPU needs its driver and the CUDA compiler, the disk and the memory need room, and Hugging Face must answer.
\item The engine: a clone of the repository \code{--repo} names into \code{~/TandemLLM/src} at the pinned release, \code{v0.13.0}.
\item A Python venv in \code{~/TandemLLM/venv} with the pinned packages; torch and torchvision come from the CUDA 13.0 index. If neither \code{uv} nor Python 3.11 is on the machine, the installer fetches \code{uv} first.
\item The weights, into the normal Hugging Face cache: the vendor's FP8 checkpoint (it supplies the tensors the NVFP4 overlays do not replace, and the tokenizer), the four NVFP4 files and both fine-tuned drafters. With \code{--setup}, only Wren's model; \code{use-profile.sh} fetches the rest later.
\item A lookup corpus built from public text (one English and one German Wikipedia shard, plus Python sources from the venv's own packages, about 31 million tokens). \code{--no-corpus} skips it.
\item The configuration: \code{~/TandemLLM/run/ops/serve.env}, the release's \code{ops/serve.env}\index{ops/serve.env@\texttt{ops/serve.env}} with this machine's paths, and a control script, \code{~/TandemLLM/bin/tandem}.
\item The start: the engine starts, the script waits for \code{/health}, and a smoke test prints tokens a second. \code{--no-start} stops after step 6.
\end{enumerate}

The weights are public. \code{--hf-token} (or \code{HF_TOKEN}) is needed only for a repository that is private or gated, and a stored \code{hf auth login} works too. Table~\ref{tab:ref-install-flags} lists every option; the most used ones are \code{--setup} (Section~\ref{sec:new-box}), \code{--repo} (another repository, such as a fork), \code{--dir} (default \code{~/TandemLLM}), \code{--port} and \code{--host} (default \code{127.0.0.1:8000}), \code{--ref} (any git ref; \code{v0.2.0-staircut} is the paper's release), \code{--hf-token}, \code{--systemd} for a systemd user unit, \code{--dry-run} to print every step and change nothing, and \code{--uninstall}.

% Generated by manual/tools/gen_reference.py from the code. Do not edit by hand.

\begin{reftabletwo}{install.sh options}{tab:ref-install-flags}{\hd{Option} & \hd{What it does}}
\texttt{-{}-dir DIR} & install directory (default \textasciitilde{}/TandemLLM) \\
\texttt{-{}-port PORT} & the API port (default 8000) \\
\texttt{-{}-host ADDR} & the address to bind (default 127.0.0.1; 0.0.0.0 serves the LAN, with no auth on the OpenAI routes) \\
\texttt{-{}-ref REF} & the engine release: tag, branch or commit (default v0.13.0) \\
\texttt{-{}-repo URL} & the git repository to clone (default the public repository) \\
\texttt{-{}-hf-\allowbreak{}token TOKEN} & a Hugging Face token (or set HF\_TOKEN; a stored 'hf auth login' works too) \\
\texttt{-{}-no-\allowbreak{}corpus} & skip the lookup corpus (the lookup drafter then copies from the conversation only) \\
\texttt{-{}-no-\allowbreak{}start} & install and configure, but do not start the engine \\
\texttt{-{}-setup} & set up in the browser instead: install the code and Wren's model only, then start setup mode (the dashboard's \#/setup page; docs/setup.md); 127.0.0.1 unless -\kern0pt-host says otherwise \\
\texttt{-{}-force-\allowbreak{}setup} & with -\kern0pt-setup, on a directory that already has a configured engine \\
\texttt{-{}-systemd} & also install a systemd user unit, tandemllm.service (off by default) \\
\texttt{-{}-yes,\allowbreak{} -\allowbreak{}y} & non-interactive: answer yes to every question \\
\texttt{-{}-dry-\allowbreak{}run} & print every step and command, change nothing \\
\texttt{-{}-uninstall} & stop the engine, remove the unit and DIR (the weights stay in the HF cache) \\
\texttt{-h,\allowbreak{} -{}-help} & this text \\
\end{reftabletwo}
\noindent Environment variables the installer reads: \texttt{TANDEM\_\allowbreak{}REF}, \texttt{TANDEM\_\allowbreak{}REPO}, \texttt{TANDEM\_\allowbreak{}DIR}, \texttt{TANDEM\_\allowbreak{}PORT}, \texttt{TANDEM\_\allowbreak{}HOST}, \texttt{HF\_\allowbreak{}TOKEN}, \texttt{HF\_\allowbreak{}HOME}.


The installer's corpus is not the corpus the paper's numbers come from; that one is not published. With the installer's corpus, or with \code{--no-corpus}, the lookup drafter proposes different continuations, so speeds can differ from the paper's. How much is not measured. On a CUDA GPU other than the GB10 the installer turns off the skinny NVFP4 kernel, which is built for \code{sm_121a} only, and StairCut's cost tables still come from the Spark, so such GPUs are best effort.

\section{After the install}

The install directory holds everything the installer wrote:

\begin{table}[htbp]
\centering\small
\caption{The layout of an installed engine.}\label{tab:install-layout}
\begin{tabularx}{\linewidth}{@{}L{.32\linewidth}Y@{}}
\toprule
\hd{path} & \hd{what} \\
\midrule
\code{~/TandemLLM/src} & the repository at the installed release \\
\code{~/TandemLLM/venv} & the Python environment, about 8~GiB with the CUDA wheels \\
\code{~/TandemLLM/corpus} & the lookup corpus \\
\code{~/TandemLLM/run/ops/serve.env} & the served profile with this machine's paths, rewritten on every run of the installer \\
\code{~/TandemLLM/run/local.env} & your own settings; \code{serve.env} reads it last and the installer never touches it \\
\code{~/TandemLLM/state} & the tokens file (the metrics token, and the admin password if you set one), once you create it \\
\code{~/TandemLLM/bin/tandem} & the control script \\
\code{~/.qwen38-spark-engine/} & the engine's state: the usage ledger, the conductor's documents and journal, memory notes, the disk tier, small models \\
\bottomrule
\end{tabularx}
\end{table}

The control script takes one command:

\begin{shell}
~/TandemLLM/bin/tandem start      # ops/start.sh with run/ops/serve.env; waits for /health
~/TandemLLM/bin/tandem stop 120   # SIGTERM, wait for the drain, kill only after it
~/TandemLLM/bin/tandem restart
~/TandemLLM/bin/tandem status
~/TandemLLM/bin/tandem smoke      # /v1/models and one short completion, with tok/s
~/TandemLLM/bin/tandem logs       # follow the engine's log
\end{shell}

\section{The first request}

The server speaks the OpenAI API on port 8000. With the conductor on, as in the shipped profile, it serves one model called \code{TandemLLM}, and any model name a client sends is served as \code{TandemLLM}:

\begin{shell}
curl -s localhost:8000/v1/chat/completions -H 'Content-Type: application/json' \
    -d '{"model": "TandemLLM", "messages": [{"role": "user", "content": "Say hello."}]}'
\end{shell}

Point other clients at the same address:

\begin{table}[htbp]
\centering\small
\caption{Pointing a client at the engine.}\label{tab:clients}
\begin{tabularx}{\linewidth}{@{}L{.2\linewidth}Y@{}}
\toprule
\hd{client} & \hd{setting} \\
\midrule
an OpenAI SDK & \code{OPENAI_BASE_URL=http://<host>:8000/v1} and any key \\
opencode & an OpenAI-compatible provider with that base URL and the model \code{TandemLLM}; it gets the \code{coding} profile by itself \\
Open WebUI & an OpenAI connection to that base URL, with the connection header \code{X-OpenWebUI-Chat-Id: {{CHAT_ID}}} so each chat is its own conductor session (Section~\ref{sec:router-session}); it gets the \code{chat} profile \\
\code{curl}, a script & as above; it gets the \code{general} profile \\
\bottomrule
\end{tabularx}
\end{table}

For the engine without the conductor, send the header \code{X-Tandem: raw}\index{X-Tandem header}. That is the release's raw path, the one the benchmarks and the paper measure. Chapter~\ref{ch:conductor} explains the profiles.

\screenshot[trim=0 640 0 0, clip]{client-open-webui}{Open WebUI with \texttt{TandemLLM} as its model: a demo request on the test board in October 2026, served under the \texttt{chat} profile, with the reasoning folded into \enquote{Thought for 1 seconds}. Only the message pane is shown.}{fig:shot-open-webui}

\screenshot{client-opencode}{opencode with \texttt{TandemLLM} as its provider: the end of a demo task on the test board in October 2026, under the \texttt{coding} profile. The model added a function, tested it, asked Laya for a verdict on the test log, removed its scratch files and reported what it did. The screenshot is from a build that sent the reasoning to opencode as text between \texttt{<think>} tags; every client but the Hermes agent now gets \texttt{reasoning\_content} by default (\texttt{QWEN38\_CLIENT\_REASONING}), which opencode shows as folded reasoning parts.}{fig:shot-opencode}

The dashboard is at \code{http://<host>:8000/dashboard/}. It asks for no password; read Section~\ref{sec:access} before anyone else can reach the port.

\section{Access and sign-in}\label{sec:access}\label{sec:dash-login}\index{dashboard login}\index{admin password}\index{sign-in}

TandemLLM has no sign-in turned on by default. The engine's OpenAI routes are always open, and so is everything else until you set an admin password: the dashboard, the conductor's settings, the memory notes, the transcripts, the decision models, the configuration and the cache routes. Anyone who reaches the port can read and change them. The installer binds the engine to \code{127.0.0.1}, so only the machine itself reaches it. With \code{--host 0.0.0.0}, or when you start \code{ops/serve.env} from a checkout (its \code{HOST} is \code{0.0.0.0}), the whole local network does.

Table~\ref{tab:access} has the switches. Only \code{/metrics} has a token of its own, for a Prometheus scraper, and \code{ops/make-secrets.sh} writes that one alone unless \code{--admin} asks for the password too.

\begin{table}[htbp]
\centering\small
\caption{Who can reach what, by default and with the switches set.}\label{tab:access}
\begin{tabularx}{\linewidth}{@{}L{.27\linewidth}L{.27\linewidth}Y@{}}
\toprule
\hd{what} & \hd{by default} & \hd{with the switch} \\
\midrule
the engine's OpenAI routes & anyone who reaches the port & no switch on the engine's own port; behind the proxy, a user key from the first key on (Section~\ref{sec:users}) \\
the dashboard's reads & anyone & \code{--dashboard-login on} and a password: the token or a session \\
admin routes: conductor writes, memory, decision models, configuration, caches, the full \code{/health} & anyone & \code{QSE_ADMIN_TOKEN} (\code{ops/make-secrets.sh --admin}): the token or a session, or a tool on the machine itself \\
\code{/metrics} & anyone & \code{QSE_METRICS_TOKEN}: that token \\
the standalone proxy & binds \code{127.0.0.1}; open & \code{TANDEM_ADMIN_TOKEN}, user keys \\
\bottomrule
\end{tabularx}
\end{table}

Setting a password has a cost in this release: the dashboard has no sign-in screen. Its read-only views keep working while the login is off, but its edits and the Memory screen show the server's 401 and a banner. Remove the line from \code{secrets.env} and restart to open it again. A request from the machine itself (loopback, without proxy headers) needs no token, because the watchdog and the measurement tools rely on that.

TandemLLM serves one user. Every client that reaches the engine is that user, so a \enquote{remember that} from any client lands in the same memory notes. A request that must not touch them sends \code{X-Tandem-Memory: off}. Until sign-in and several users arrive, keep the port on loopback or a trusted network. To reach it from elsewhere, put a reverse proxy with its own sign-in, or a VPN, in front of it.

\section{The small models}\label{sec:small-models}

Three small models can run beside the main model. They are off in the shipped profile, because each needs its weights on the disk. Each is one line in \code{local.env} and a restart.

\begin{table}[htbp]
\centering\small
\caption{The small models and the line that turns each on.}\label{tab:optional-parts}
\begin{tabularx}{\linewidth}{@{}L{.22\linewidth}L{.34\linewidth}Y@{}}
\toprule
\hd{part} & \hd{set in \code{local.env}} & \hd{the model} \\
\midrule
Laya, the decision models & \code{QWEN38_DECIDER_MODELS=laya-english,laya-multilingual,laya-typed-decisions} & \code{convaiinnovations/laya} under \code{~/.qwen38-spark-engine/deciders/laya} \\
the narrator's voice & \code{QWEN38_NARRATOR_MODEL=qwen3.5-2b} & \code{Qwen/Qwen3.5-2B} under \code{~/.qwen38-spark-engine/models/qwen3.5-2b} \\
the shadow reader & \code{QWEN38_SHADOW_MODEL=qwen3.5-0.8b} & \code{Qwen/Qwen3.5-0.8B} under \code{~/.qwen38-spark-engine/models/qwen3.5-0.8b} \\
\bottomrule
\end{tabularx}
\end{table}

Download a model with \code{hf download <repo> --local-dir <dir>}. The small models load after the engine, one at a time, through the residency planner\index{residency planner}, which refuses a load that would leave the board short of memory (Section~\ref{sec:residency}). The dashboard's System view shows what is loaded.

\section{Other engines and cloud APIs beside the engine}\label{sec:engine-backends}

The engine's own server can be the one address for everything. Point \code{QSE_BACKENDS}\index{qsebackends@\texttt{QSE\_BACKENDS}} in \code{local.env} at a config database and restart:

\begin{shell}
echo 'QSE_BACKENDS=$HOME/.tandemllm/config.sqlite3' >> ~/TandemLLM/run/local.env
~/TandemLLM/bin/tandem restart
\end{shell}

Then add the backends on the dashboard's Backends screen, or as \code{TANDEM_*} variables (Chapter~\ref{ch:other-engines}). The engine keeps serving \code{TandemLLM} as before, and a request for another model goes to its backend on the handler's own thread: it never takes the engine's lock or a place in its queue, so a slow cloud answer cannot hold up the engine. The engine's conductor can run in front of each backend (\code{harness.mode}, Chapter~\ref{ch:harness-proxy}), and one ledger and the Parallel view cover them all. On a DGX Spark the other backends are cloud APIs or engines on other machines; a second GPU engine on the same Spark is not supported (Table~\ref{tab:setups}). No live run of this in-process relay is recorded yet; the standalone proxy in front of the engine is the setup that ran live.

The other way round works too: \code{tandem-dash} on a laptop or a small server, with the engine as a backend of kind \code{tandemllm} (Chapter~\ref{ch:install-dash}).

\section{Running from a checkout}

The installer does nothing you cannot do by hand. In an environment with the packages above, download the published weights and drafters with \code{hf}\index{Hugging Face}, the Hugging Face CLI (\code{pip install -U huggingface_hub}):

\begin{shell}
hf download Qwen/Qwen3.8-27B-FP8
hf download 0xBakeer/TandemLLM-Qwen3.8-27B-NVFP4 --local-dir ~/tandem/nvfp4
hf download 0xBakeer/TandemLLM-Qwen3.8-27B-DFlash2-b8 --local-dir ~/tandem/ft-b8
hf download 0xBakeer/TandemLLM-Qwen3.8-27B-DFlash2-b16 --local-dir ~/tandem/ft-b16
\end{shell}

Then edit the path lines at the top of \code{ops/serve.env}: \code{REPO} (the checkout), \code{PY} (a Python with the packages), \code{NV} (the three NVFP4 files, \code{mlp}, \code{gdn} and \code{attn}, comma-separated), \code{HEAD} (\code{head-fp8.safetensors}), \code{CKPT8}, \code{CKPT16} and \code{CORPUS}. Start it with \code{bash ops/start.sh}. That is the served profile.

In a PyTorch container that already has torch 2.13.0+cu130, add the two image packages:

\begin{shell}
pip install pillow==11.3.0 torchvision==0.28.0+cu130 \
    --extra-index-url https://download.pytorch.org/whl/cu130
\end{shell}

A server started directly with \code{python server/app.py} and no further flags is the raw engine. The conductor, the disk tier and the usage ledger come only with \code{--conductor-config}\index{conductor-config@\texttt{-\kern0pt-conductor-config}}, \code{--disk-cache}\index{disk-cache@\texttt{-\kern0pt-disk-cache}} and \code{--usage-ledger}\index{usage-ledger@\texttt{-\kern0pt-usage-ledger}}, which \code{ops/start.sh}\index{ops/start.sh@\texttt{ops/start.sh}} passes from \code{QSE_CONDUCTOR_CONFIG}\index{qseconductorconfig@\texttt{QSE\_CONDUCTOR\_CONFIG}}, \code{QSE_DISK_CACHE}\index{qsediskcache@\texttt{QSE\_DISK\_CACHE}} and \code{QSE_USAGE_LEDGER}\index{qseusageledger@\texttt{QSE\_USAGE\_LEDGER}}. To try the engine with the vendor's base drafter instead of the fine-tuned ones:

\begin{shell}
hf download z-lab/Qwen3.8-27B-DFlash2
python server/app.py --port 8000 --max-len 32768 --drafter dflash2 --dflash2-path greedy \
    --nvfp4 ~/tandem/nvfp4/mlp.safetensors,~/tandem/nvfp4/gdn.safetensors,~/tandem/nvfp4/attn.safetensors \
    --fp8-head ~/tandem/nvfp4/head-fp8.safetensors
\end{shell}

\section{Building the weights yourself}

The published NVFP4 set was quantised once from Qwen's BF16 release with GPTQ on the NVFP4 grid. Rebuilding it needs the BF16 checkpoint and a GPU with about 80~GB of memory at the default \code{--h-budget-gb 22}; it took 25 minutes on an RTX PRO 6000 (\measured):

\begin{shell}
hf download Qwen/Qwen3.8-27B
python tools/calib_corpus.py --out ~/calib
python tools/quant_nvfp4.py build --model <the BF16 snapshot> --methods gptq --out-dir ~/nvfp4-bf16 \
    --corpus ~/calib/calib-code.txt:80000,~/calib/calib-en.txt:60000,~/calib/calib-de.txt:16000
python tools/quant_head.py build --out ~/nvfp4-bf16/head-fp8.safetensors --ratios 1.0,0.95,0.90
\end{shell}

The build writes \code{gptq/mlp.safetensors}, \code{gdn.safetensors} and \code{attn.safetensors}; \code{tools/quality_gate.py --plan} scores them against the BF16 model. The engine loads them over the FP8 checkpoint, as it does the published files. Chapter~\ref{ch:quantisation} explains the method and the gate.

\section{Profiles}\index{profile (serve.env)}

A serving profile is a \code{serve.env} file: which weights, which drafters, which settings. Three ship with the NVFP4 set. They differ only in their weights; drafters, StairCut, caches and the exactness gate are the same in all three. The weight profiles under \code{ops/profiles/} are the other kind: \code{nvfp4}, six \code{exl3-<bpw>} sizes of the 27B and five \code{flashnext-<bpw>} packs, which \code{ops/use-profile.sh} switches between one engine at a time (Section~\ref{sec:two-models}).

\begin{table}[htbp]
\centering\small
\caption{The three serving profiles. Only the NVFP4 profile's speed is measured on this release; the bytes a step are \arith{} from the tensor sizes.}\label{tab:profiles-env}
\begin{tabularx}{\linewidth}{@{}L{.17\linewidth}L{.22\linewidth}Yr@{}}
\toprule
\hd{profile} & \hd{file} & \hd{weights} & \hd{bytes a step} \\
\midrule
NVFP4 (served) & \code{ops/serve.env} & the published set: every projection NVFP4, FP8 head & 15.0 GB \\
balanced & \code{ops/serve-balanced.env} & the published NVFP4 MLPs and FP8 head; the checkpoint's FP8 recurrent and attention layers & about 18.2 GB \\
FP8 & \code{ops/serve-fp8.env} & the checkpoint's own FP8 projections and BF16 head & 26.9 GB \\
\bottomrule
\end{tabularx}
\end{table}

The balanced and FP8 profiles source \code{ops/serve.env} and change two lines, \code{NV} and \code{HEAD}. They inherit StairCut's verify prices, which were measured on the NVFP4 profile. The router corrects their level online, but the staircase's shape for those two profiles is not measured, and neither is their speed. Do not confuse these files with the conductor's profiles (\code{coding}, \code{chat} and the others), which are documents the dashboard edits.

\section{Upgrading and uninstalling}\index{upgrade}\index{uninstall}

To upgrade, run the installer again with the same \code{--dir}. It checks out the new release, installs the new lock and rewrites \code{serve.env} from the release's template; \code{local.env} stays. The usage ledger migrates in place at the first start: columns are added and every row is kept. A box installed with \code{--setup} and switched to a profile since then upgrades with \code{--setup --force-setup} and the profile switch after it, as Section~\ref{sec:update-setup} shows.

Coming from 0.3.0, a client sees these changes:

\begin{itemize}
\item The conductor is on in the shipped profile. \code{/v1/models} lists \code{TandemLLM} alone, every response carries the headers \code{X-Tandem}\index{x-tandem@\texttt{X-Tandem}} and \code{X-Tandem-Profile}\index{x-tandem-profile@\texttt{X-Tandem-Profile}}, and every request runs through the router, with a real cap on the thinking per effort level. \code{X-Tandem: raw} gives the release path of 0.3.0. A benchmark client that cannot send a header should run against a server with \code{CONDUCTOR=raw} and an empty \code{QSE_CONDUCTOR_CONFIG} in \code{local.env}.
\item The disk tier and the session spill are on: up to 50~GB under \code{~/.qwen38-spark-engine/cache/disk}, never below 50~GB free on the disk. \code{QSE_DISK_CACHE=off} turns both off.
\item A streamed request that finds the engine busy waits as long as a request may run, with an SSE comment every 10~s, instead of a 429 after 240~s. A short request is served between two blocks of the running turn.
\item The image processor's backend is pinned to torchvision.
\item A malformed \code{tools} field is a 400 that names the item.
\item The memory limits are enforced: 70~GiB between requests, 77~GiB during a prefill.
\end{itemize}

To run 0.4.0 as 0.3.0 was served, put \code{CONDUCTOR=raw}, an empty \code{QSE_CONDUCTOR_CONFIG=} and \code{QSE_DISK_CACHE=off} in \code{local.env}.

To roll back, run the installer with \code{--ref} set to the earlier tag; the ledger keeps its added columns and the older release ignores them. To uninstall:

\begin{shell}
bash ~/TandemLLM/src/install.sh --uninstall --yes
\end{shell}

This stops the engine, removes the systemd unit if one was installed, and deletes the install directory. It keeps the weights in the Hugging Face cache, which other tools may share, and \code{~/.qwen38-spark-engine/} with the conductor's documents, the journal, memory notes, transcripts, the disk tier and the small models. Delete that directory to remove them too.
